\NSPage{013}%
The correction cycle treats the remaining linear and quadratic interactions, including the
curl and cutoff errors in \NSi{NS-F-P013-001}{\mathscr R(u_B+w,p_B+\pi)}.

\begin{figure}[htbp]
  \centering
  \NSFigurePlaceholder{NS-FIG-5}
  \caption{Construction of the background and the first oscillatory correction.
  The averaged pulse flux cancels the leading background stress divergence.
  The remaining residual is corrected by the cycle in Figure~\ref{fig:6}. A flat remainder
  and all curl and cutoff terms are retained.}
  \label{fig:5}
\end{figure}

\subsection{Correction and summation}\label{sec:3.4}
The averaged wave flux is chosen to cancel the leading stress divergence associated with
\NSi{NS-F-P013-002}{T} (Propositions~\ref{prop:7.5} and~\ref{prop:9.5}). The remaining angular
harmonics, auxiliary means, and radial integral defects require separate corrections. We
treat them in a fixed cycle and recompute the full residual after each operation, as in
Proposition~\ref{prop:9.6}, so newly created terms enter the next stage.

Starting from \NSi{NS-F-P013-003}{(u_B+w,p_B+\pi)}, we reconstruct the pressure and correct the angularly
averaged velocity and its radial integral defects. Proposition~\ref{prop:9.5} proves that the resulting
pair \NSi{NS-F-P013-004}{(u^{[0]},p^{[0]})} satisfies the bounds required for the correction cycle. Square brackets count
these cycles; the superscript \NSi{NS-F-P013-005}{(0)} continues to denote the leading field
\NSi{NS-F-P013-006}{(u^{(0)},p^{(0)})}.

At stage \NSi{NS-F-P013-007}{j}, we add a divergence-free velocity increment
\NSi{NS-F-P013-008}{\delta u_j} and a pressure increment \NSi{NS-F-P013-009}{\delta p_j}:
\NSFormula{NS-F-P013-010}{display}{none}
\[
u^{[j+1]}=u^{[j]}+\delta u_j,
\qquad
p^{[j+1]}=p^{[j]}+\delta p_j.
\]
The residual changes according to
\NSFormula{NS-F-P013-011}{display}{none}
\[
\begin{aligned}
\mathscr R(u^{[j+1]},p^{[j+1]})
  &=\mathscr R(u^{[j]},p^{[j]})+\mathcal L_{u^{[j]}}(\delta u_j,\delta p_j)\\
  &\qquad{}+\nabla\mathbin{\cdot}(\delta u_j\otimes\delta u_j).
\end{aligned}
\]

\NSPage{014}%
Here \NSi{NS-F-P014-001}{\mathcal L_v} is the linear residual operator defined above, with background
\NSi{NS-F-P014-002}{v}. Thus canceling a selected source also introduces linear remainders and quadratic interactions. The four
operations in Proposition~\ref{prop:9.6} control these terms as follows.
\begin{enumerate}[label=\arabic*.,leftmargin=*]
\item We solve the inhomogeneous wave-amplitude equations for the supported nonzero
angular Fourier modes. Their principal operator cancels the prescribed source. The
remaining linear terms and the interactions of the new pulses enter the next residual
(Proposition~\ref{prop:9.1} and Lemma~\ref{lem:9.2}).
\item We construct signed amplitude increments using the fixed positive leading amplitudes.
Their symmetrized cross covariance with the leading pulses supplies the prescribed
correction to the averaged stress. Interactions with earlier corrections and the quadratic
self-interaction remain as higher-order terms (Proposition~\ref{prop:7.6}, \eqref{eq:9.13}).
\item We correct the angularly averaged residual with zero auxiliary average by inverting the
fast auxiliary-time derivative \eqref{eq:8.20}. The axial increment is realized through a vector
potential in \eqref{eq:8.14}, which also supplies its radial velocity. The reconstruction remainder,
whose Cartesian space-time derivatives vanish faster than every power of
\NSi{NS-F-P014-003}{q}, is retained in the residual.
\item We solve the five radial moment equations in \eqref{eq:8.25}. Two equations preserve the zero
angular-momentum and axial-flux integrals in \eqref{eq:9.10}. Three cancel the linear contributions
to \NSi{NS-F-P014-004}{(\mathcal P,\mathcal J_\theta,\mathcal J_z)}, the radial integral defects for pressure and tangential momentum defined
in \eqref{eq:8.12} and \eqref{eq:8.15}; their nonlinear remainders have improved decay \eqref{eq:8.27}. The radial
integral constructions keep the velocity and pressure corrections compactly supported.
\end{enumerate}
Pressure is reconstructed from the updated radial equation after each operation, and its contribution
to the axial equation is included in the new residual. All products use the updated
fields, including the curl and cutoff corrections. The background, leading amplitudes, and
inverse operators remain fixed throughout the induction.

The parameter \NSi{NS-F-P014-005}{\sigma_j} in \eqref{eq:9.8} measures residual decay. For Cartesian derivatives of order
\NSi{NS-F-P014-006}{m}, the bound in \eqref{eq:9.18} contains the power
\NSi{NS-F-P014-007}{q^{h\sigma_j-K_m}}, where \NSi{NS-F-P014-008}{K_m} is independent of
\NSi{NS-F-P014-009}{j}, together with logarithmic factors and the separately retained remainder. Constants and logarithmic powers
may depend on \NSi{NS-F-P014-010}{j,m}. Initialization and one correction cycle give
\NSFormula{NS-F-P014-011}{display}{none}
\[
\sigma_0=\frac15,
\qquad
\sigma_{j+1}=\sigma_j+\frac1{10},
\qquad
\sigma_j=\frac15+\frac j{10}\longrightarrow\infty.
\]
The cumulative velocity bounds \eqref{eq:9.9}, support conditions, and two preserved integrals \eqref{eq:9.10}
persist, so the cycle can be repeated. All finite stages are defined on the common domain of
Lemma~\ref{lem:9.7}.

To sum the corrections, we apply cutoffs to their vector potentials and pressures, then take
curls; the direct axisymmetric azimuthal increments receive the same cutoffs. Each cutoff
equals one sufficiently close to \NSi{NS-F-P014-012}{q=0}, and its support shrinks with the stage. The bounds in
\eqref{eq:9.17} have a loss in the power of \NSi{NS-F-P014-013}{q} that depends on derivative order but is independent of
the stage. They allow the cutoffs to be chosen so that every differentiated tail satisfies \eqref{eq:5.35}.
The sum is locally finite for \NSi{NS-F-P014-014}{q>0} and defines smooth fields
\NSi{NS-F-P014-015}{(u_{\mathrm{loc}},p_{\mathrm{loc}})} for \NSi{NS-F-P014-016}{t<1}. Comparing
their residual with that of a fixed finite stage proves \eqref{eq:9.20}, while preserving the leading
velocity growth (Proposition~\ref{prop:9.9}).

The summation gives the required local velocity growth. To pass to the whole space, we
also need a vector-potential representation and regularity up to time one away from the
origin. Proposition~\ref{prop:9.9} proves these properties together with flatness of the residual. In the
theorem, we write \NSi{NS-F-P014-017}{(u,p)} for the local pair
\NSi{NS-F-P014-018}{(u_{\mathrm{loc}},p_{\mathrm{loc}})}.

\NSPage{015}%
\begin{figure}[htbp]
  \centering
  \NSFigurePlaceholder{NS-FIG-6}
  \caption{Residual correction and completion, continuing Figure~\ref{fig:5}. The
  correction cycle improves the residual exponent \NSi{NS-F-P015-001}{\sigma_j} while retaining all linear
  and nonlinear terms. The flatness estimate holds uniformly for
  \NSi{NS-F-P015-002}{0\leq X\leq X_1}
  as \NSi{NS-F-P015-003}{q\downarrow0}, for each finite \NSi{NS-F-P015-004}{X_1}. Incompressibility is preserved throughout.}
  \label{fig:6}
\end{figure}

We write \NSi{NS-F-P015-005}{\partial_x^\alpha\partial_t^b} for Cartesian space-time derivatives, where
\NSi{NS-F-P015-006}{\alpha} is a spatial multi-index and
\NSi{NS-F-P015-007}{b\geq0} is an integer.

\begin{theorem}\label{thm:3.1}
There are constants \NSi{NS-F-P015-008}{0<h<1/100},
\NSi{NS-F-P015-009}{q_*>0}, \NSi{NS-F-P015-010}{0<X_a<X_{\mathrm{ext}}<\infty}, and smooth fields on
\NSFormula{NS-F-P015-011}{display}{none}
\[
\Omega_* = \{(x,t):\tau>0,\ q<q_*\}
\]
with the following properties.
\begin{enumerate}[label=\textup{(\roman*)},leftmargin=*]
\item There is a vector potential \NSi{NS-F-P015-012}{\mathcal A} and an axisymmetric scalar
\NSi{NS-F-P015-013}{\mathcal B} such that
\NSFormula{NS-F-P015-014}{display}{3.3}
\begin{equation}
u=\operatorname{curl}\mathcal A+\mathcal B e_\theta.
\tag{3.3}\label{eq:3.3}
\end{equation}
Both \NSi{NS-F-P015-015}{\mathcal A} and \NSi{NS-F-P015-016}{\mathcal B e_\theta} are smooth in Cartesian coordinates at the axis. The field is exactly divergence-free.

\item Every Cartesian space-time derivative of
\NSi{NS-F-P015-017}{\mathcal A,\mathcal B e_\theta,p} is uniformly bounded on compact spatial
subsets with \NSi{NS-F-P015-018}{c\leq q\leq c'<q_*}, for each
\NSi{NS-F-P015-019}{0<c<c'}, up to the one-sided boundary
\NSi{NS-F-P015-020}{\tau=0}. The resulting limits of the derivatives are compatible under spatial and time differentiation.

\item For every \NSi{NS-F-P015-021}{\alpha,b,N} and every finite
\NSi{NS-F-P015-022}{X_1},
\NSFormula{NS-F-P015-023}{display}{3.4}
\begin{equation}
\lvert\partial_x^\alpha\partial_t^b\mathscr R(u,p)\rvert
\leq C_{\alpha,b,N,X_1}q^N
\qquad (0\leq X\leq X_1,\ q\downarrow0).
\tag{3.4}\label{eq:3.4}
\end{equation}
For \NSi{NS-F-P015-024}{X\geq X_{\mathrm{ext}}} the residual is identically zero, and
\NSFormula{NS-F-P015-025}{display}{3.5}
\begin{equation}
\mathcal A=0,
\qquad \mathcal B=K(r,\tau),
\qquad p=-\int_r^\infty\frac{K(\rho,\tau)^2}{\rho}\,d\rho,
\qquad K(r,\tau)=r^{-1-2h}\mathcal H_{\mathrm{ext}}(\tau/r^2).
\tag{3.5}\label{eq:3.5}
\end{equation}
Here \NSi{NS-F-P015-026}{\mathcal H_{\mathrm{ext}}} is smooth on
\NSi{NS-F-P015-027}{[0,\infty)}, with each fixed derivative bounded. The radial swirl heat equation satisfied by
\NSi{NS-F-P015-028}{K} is
\NSFormula{NS-F-P015-029}{display}{none}
\[
-\partial_\tau K=\partial_r^2K+r^{-1}\partial_rK-r^{-2}K.
\]

\NSPage{016}%
In this statement the second argument of \NSi{NS-F-P016-001}{K} is
\NSi{NS-F-P016-002}{\tau}; when written as \NSi{NS-F-P016-003}{K(r,t)} below, the same field
is evaluated at \NSi{NS-F-P016-004}{\tau=1-t}.

\item For some fixed \NSi{NS-F-P016-005}{X_{\mathrm{in}}\in(0,X_a)} and
\NSi{NS-F-P016-006}{e_0>0},
\NSFormula{NS-F-P016-007}{display}{3.6}
\begin{equation}
u_\theta(\sqrt{2X_{\mathrm{in}}\tau},0,0,1-\tau)
=\tau^{-A}\bigl(e_0+O(\tau^{2h})\bigr)
\qquad(\tau\downarrow0).
\tag{3.6}\label{eq:3.6}
\end{equation}
The entries on the left specify cylindrical position and time.
\end{enumerate}
\end{theorem}

\subsection{Localization and completion}\label{sec:3.5}
Finally, using the decomposition in \eqref{eq:3.3}, we multiply the
vector potential \NSi{NS-F-P016-008}{\mathcal A}, azimuthal coefficient
\NSi{NS-F-P016-009}{\mathcal B}, and pressure \NSi{NS-F-P016-010}{p_{\mathrm{loc}}} by
\NSi{NS-F-P016-011}{c=\chi_x\chi_t}. Here \NSi{NS-F-P016-012}{\chi_x} is smooth
and axisymmetric with compact support \NSi{NS-F-P016-013}{\mathcal K}, and
\NSi{NS-F-P016-014}{\chi_t} is smooth with \NSi{NS-F-P016-015}{\chi_t=0} on
\NSi{NS-F-P016-016}{[0,1-\tau_0]} for some \NSi{NS-F-P016-017}{0<\tau_0<1/2}. Both cutoffs equal one near
\NSi{NS-F-P016-018}{(0,1)}. Since \NSi{NS-F-P016-019}{q\asymp\tau+|z|^{1/D}}, we can
choose their supports so that \NSi{NS-F-P016-020}{q<q_*/2} wherever
\NSi{NS-F-P016-021}{c\neq0} and \NSi{NS-F-P016-022}{t<1}. We take the curl after
multiplication and extend by zero outside the cutoff support:
\NSFormula{NS-F-P016-023}{display}{none}
\[
u=\operatorname{curl}(c\mathcal A)+c\mathcal B e_\theta,
\qquad p=cp_{\mathrm{loc}},
\qquad
u=cu_{\mathrm{loc}}+\nabla c\times\mathcal A.
\]
The curl term and the axisymmetric azimuthal term are separately divergence-free. Cartesian
smoothness of the chosen representatives and the margin from \NSi{NS-F-P016-024}{q=q_*} give a smooth zero
extension. Thus \NSi{NS-F-P016-025}{u,p} have fixed spatial support in
\NSi{NS-F-P016-026}{\mathcal K}, vanish on the stated initial time interval,
and agree with the local fields near \NSi{NS-F-P016-027}{(0,1)} (Proposition~\ref{prop:10.1}).

For \NSi{NS-F-P016-028}{t<1}, we set \NSi{NS-F-P016-029}{f=\mathscr R(u,p)}. In
\NSi{NS-F-P016-030}{f-c\mathscr R(u_{\mathrm{loc}},p_{\mathrm{loc}})}, differentiating the cutoffs produces
terms such as \NSi{NS-F-P016-031}{(\partial_t c-\Delta c)u_{\mathrm{loc}}} and
\NSi{NS-F-P016-032}{p_{\mathrm{loc}}\nabla c}; the curl correction
\NSi{NS-F-P016-033}{\nabla c\times\mathcal A} contributes its derivatives
and advection products. The change in self-advection also contains
\NSi{NS-F-P016-034}{(c^2-c)(u_{\mathrm{loc}}\mathbin{\cdot}\nabla)u_{\mathrm{loc}}}.
These additional terms are supported in the cutoff transition regions, away from a neighborhood
of \NSi{NS-F-P016-035}{(0,1)}.

Near \NSi{NS-F-P016-036}{(0,1)}, the cutoffs equal one, so
\NSi{NS-F-P016-037}{f} is the local residual. Its flatness estimate \eqref{eq:3.4} on
\NSi{NS-F-P016-038}{X\leq X_{\mathrm{ext}}}, together with the identically zero residual for larger
\NSi{NS-F-P016-039}{X}, gives the bound \eqref{eq:10.9}
for every Cartesian space-time derivative. On the cutoff transition regions near time one,
Theorem~\ref{thm:3.1}\textup{(ii)} bounds all derivatives of the actual potentials and pressure where
\NSi{NS-F-P016-040}{q} stays positive. In the remaining region,
\NSi{NS-F-P016-041}{r} stays positive while \NSi{NS-F-P016-042}{q\to0}; there
\NSi{NS-F-P016-043}{\mathcal A=0}, and the heat bounds \eqref{eq:10.8} and the normalized pressure integral in \eqref{eq:3.5} give the derivative limits.
These estimates, the uniform flatness bound near the origin, and the fixed support
\NSi{NS-F-P016-044}{\mathcal K} give compatible limits uniformly on
\NSi{NS-F-P016-045}{\R^3} for every derivative of \NSi{NS-F-P016-046}{f} (Lemma~\ref{lem:10.2}). Lemma~\ref{lem:10.3}
realizes these limits from \NSi{NS-F-P016-047}{t>1} with a force supported in
\NSi{NS-F-P016-048}{\mathcal K\times[0,2]}:
\NSFormula{NS-F-P016-049}{display}{none}
\[
f\in C_c^\infty(\R^3\times(0,\infty);\R^3),
\qquad
\mathscr R(u,p)=f
\qquad(0\leq t<1).
\]
The growth path in Theorem~\ref{thm:3.1}\textup{(iv)} eventually lies where the cutoffs equal one. Hence the
localized velocity retains the asymptotic \eqref{eq:3.6} and becomes unbounded as
\NSi{NS-F-P016-050}{t\uparrow1}.

The concentration scales are compatible with bounded energy. The core
\NSi{NS-F-P016-051}{\mathcal C_\tau} has volume of
order \NSi{NS-F-P016-052}{\tau^{3/2-h}}. Its leading kinetic energy
\NSi{NS-F-P016-053}{E_{\mathrm{core}}} and the integral of the squared radial derivatives
of its leading velocity, \NSi{NS-F-P016-054}{D_{\mathrm{core}}}, have scales
\NSFormula{NS-F-P016-055}{display}{none}
\[
E_{\mathrm{core}}\asymp\tau^{3/2-h}\tau^{-1-2h}=\tau^{1/2-3h},
\qquad
D_{\mathrm{core}}\asymp\tau^{3/2-h}\tau^{-2-2h}=\tau^{-1/2-3h}.
\]
The factors are the regional volume and the squared velocity or squared radial derivative.
Since \NSi{NS-F-P016-056}{h<1/6},
\NSFormula{NS-F-P016-057}{display}{none}
\[
\tau^{1/2-3h}\longrightarrow0,
\qquad
\int_0^{\tau_0}\tau^{-1/2-3h}\,d\tau
=\frac{\tau_0^{1/2-3h}}{1/2-3h}<\infty
\qquad(\tau_0>0).
\]

\NSPage{017}%
Lemma~\ref{lem:10.4} proves the global energy and integrated dissipation bounds for the localized
fields directly from their equation.

Finally, Lemma~\ref{lem:10.5} applies to any smooth solution with the same force and zero initial
datum whose kinetic energy is uniformly bounded. It must agree with
\NSi{NS-F-P017-001}{u} on every \NSi{NS-F-P017-002}{[0,T]},
\NSi{NS-F-P017-003}{T<1}. A global smooth solution would therefore agree with the constructed velocity for
\NSi{NS-F-P017-004}{0\leq t<1}, contradicting its boundedness on a compact neighborhood of
\NSi{NS-F-P017-005}{(0,1)} along the growth path. This proves Theorem~\ref{thm:1.1}.

\subsection{Notation}\label{sec:3.6}
All profile choices are fixed before \NSi{NS-F-P017-006}{q\downarrow0}. Constants may depend on those
choices and on the stated derivative order, but are independent of the concentration scale.
The notation \NSi{NS-F-P017-007}{O(q^a)} means a bound by
\NSi{NS-F-P017-008}{Cq^a} on the specified domain; uniformity in other
parameters is stated when needed. The background expansion order
\NSi{NS-F-P017-009}{n}, correction stage \NSi{NS-F-P017-010}{j},
dyadic band index \NSi{NS-F-P017-011}{\ell}, and Fourier harmonic index are distinct.

\paragraph{Averaging conventions.}
For a scalar field \NSi{NS-F-P017-012}{F(r,\theta,z,t,Y)} on the extended domain, with the
auxiliary variable \NSi{NS-F-P017-013}{Y\in\T^2=\R^2/\Z^2} independent of the physical coordinates, we define
\NSFormula{NS-F-P017-014}{display}{3.7}
\begin{equation}
\begin{aligned}
\langle F\rangle_\theta(r,z,t,Y)
  &:=\frac1{2\pi}\int_0^{2\pi}F(r,\vartheta,z,t,Y)\,d\vartheta,\\
\langle F\rangle_Y(r,\theta,z,t)
  &:=\int_{\T^2}F(r,\theta,z,t,Y')\,dY',
  \qquad \int_{\T^2}1\,dY'=1.
\end{aligned}
\tag{3.7}\label{eq:3.7}
\end{equation}
Thus angular averaging holds the independent auxiliary coordinate fixed, while auxiliary
averaging holds all physical coordinates fixed. For vector and tensor fields, these averages
act on cylindrical components. Their composition, also called the fast average for wave fields,
is
\NSFormula{NS-F-P017-015}{display}{3.8}
\begin{equation}
\langle F\rangle
:=\langle\langle F\rangle_\theta\rangle_Y
=\langle\langle F\rangle_Y\rangle_\theta
=\frac1{2\pi}\int_{\T^2}\int_0^{2\pi}F(r,\vartheta,z,t,Y)\,d\vartheta\,dY.
\tag{3.8}\label{eq:3.8}
\end{equation}
These averages act before restriction to the physical phase map \eqref{eq:6.3}. They hold
\NSi{NS-F-P017-016}{(r,z,t)} fixed,
or equivalently \NSi{NS-F-P017-017}{(R,Z,T)} within a fixed normalized chart. The normalized Haar averages
agree on the absolute, band, and common auxiliary tori by \eqref{eq:6.19}. In the wave construction,
this composition specifies the entire fast average, including its normalization. For an angular
mean coefficient, we define its auxiliary zero-mean part by
\NSFormula{NS-F-P017-018}{display}{3.9}
\begin{equation}
f^\circ:=f-\langle f\rangle_Y,
\qquad
\langle f^\circ\rangle_Y=0.
\tag{3.9}\label{eq:3.9}
\end{equation}
The profile construction also uses a separate auxiliary loop angle
\NSi{NS-F-P017-019}{\theta'} of period \NSi{NS-F-P017-020}{2\pi}: its average
is \NSi{NS-F-P017-021}{\langle f\rangle_{\theta'}=(2\pi)^{-1}\int_0^{2\pi}f(\theta')\,d\theta'}, as in Lemmas~\ref{lem:C.1} and~\ref{lem:4.11}. Bare brackets in that local profile
convention refer to this loop average. We also use the radial average
\NSFormula{NS-F-P017-022}{display}{3.10}
\begin{equation}
\mathcal A_X(f)(X,\eta)=\frac1X\int_0^X f(X',\eta)\,dX'.
\tag{3.10}\label{eq:3.10}
\end{equation}
The auxiliary torus parametrizes oscillations; the torus in Corollary~\ref{cor:10.6} is the spatial domain
of the periodic solution.

\paragraph{Normalized operators and coefficient classes.}
In a dyadic chart, \NSi{NS-F-P017-023}{Q=2^{-\ell}} is fixed,
\NSi{NS-F-P017-024}{q\asymp Q}, and
\NSFormula{NS-F-P017-025}{display}{none}
\[
(R,Z,T)=(r/\sqrt Q,z/Q^D,\tau/Q),
\qquad
\varepsilon=Q^h,
\qquad
S_*=\ell^2.
\]
The star on a spatial operator includes the factor \NSi{NS-F-P017-026}{\sqrt Q} converting physical spatial differentiation to chart units. Physical differentiation includes the chain rule for the auxiliary
phase map. Using the radial operator \NSi{NS-F-P017-027}{\mathfrak r} of \eqref{eq:6.4}, we put
\NSi{NS-F-P017-028}{D_r=\sqrt Q\,\mathfrak r},
\NSi{NS-F-P017-029}{D_z=\sqrt Q\,\partial_z=\varepsilon\partial_Z}, and
\NSPage{018}%
\NSi{NS-F-P018-001}{D_\theta=R^{-1}\partial_\theta}, as in \eqref{eq:6.6}. For a scalar
\NSi{NS-F-P018-002}{f} and a vector
\NSi{NS-F-P018-003}{a=(a_r,a_\theta,a_z)} in cylindrical components,
our conventions are
\NSFormula{NS-F-P018-004}{display}{3.11}
\begin{equation}
\begin{aligned}
\nabla_*f&=(D_rf,D_\theta f,D_zf),\\
\operatorname{div}_*a&=(D_r+R^{-1})a_r+D_\theta a_\theta+D_za_z,\\
\operatorname{curl}_*a&=(D_\theta a_z-D_za_\theta,\ D_za_r-D_ra_z,\ (D_r+R^{-1})a_\theta-D_\theta a_r).
\end{aligned}
\tag{3.11}\label{eq:3.11}
\end{equation}
The factors normalizing the fields themselves are recorded separately in \eqref{eq:8.11}. The physical
momentum residual is \NSi{NS-F-P018-005}{\mathscr R} from \eqref{eq:3.1}.

The classes \NSi{NS-F-P018-006}{\mathsf W_\alpha,\mathsf M_\alpha,\mathsf S_\alpha} record wave amplitudes, angular mean coefficients, and radial
moments in normalized units. A coefficient derivative differentiates the amplitude with
its oscillatory exponential factored out, holding the band, label, and harmonic fixed. For
each fixed correction stage and coefficient derivative order,
\NSi{NS-F-P018-007}{\partial^I} denotes ordinary derivatives
in \NSi{NS-F-P018-008}{(R,Z,T,H)}, with \NSi{NS-F-P018-009}{H} the common auxiliary coordinate of \eqref{eq:6.17}; for moments, only
\NSi{NS-F-P018-010}{(Z,T)} are differentiated. The defining bounds have the following forms. The first two hold on
\NSi{NS-F-P018-011}{X_a<X<X_b}, and the wave bound is restricted to its labelled rectangle:
\NSFormula{NS-F-P018-012}{display}{none}
\[
\begin{aligned}
f\in \mathsf M_\alpha
&\quad\Longrightarrow\quad
|\partial^If|\leq C_{j,I}\varepsilon^\alpha S_*^{b_{j,I}}\zeta\delta^{-d_{j,I}},\\
a\in \mathsf W_\alpha
&\quad\Longrightarrow\quad
|\partial^Ia|\leq C_{j,I}\varepsilon^\alpha S_*^{b_{j,I}}\sqrt\zeta\,\delta^{-d_{j,I}}P_v
\qquad(0\leq v\leq L_s),\\
F\in \mathsf S_\alpha
&\quad\Longrightarrow\quad
|\partial_{Z,T}^IF|\leq C_{j,I}\varepsilon^\alpha S_*^{b_{j,I}}.
\end{aligned}
\]
Here \NSi{NS-F-P018-013}{\zeta,\delta} are the radial-edge weights of \eqref{eq:6.23}, and
\NSi{NS-F-P018-014}{P_v} is the pulse envelope of \eqref{eq:7.12}, used
as a pointwise weight. The constants \NSi{NS-F-P018-015}{C_{j,I}>0},
\NSi{NS-F-P018-016}{b_{j,I},d_{j,I}\geq0} may differ by row and depend
on the fixed stage, derivative order, and profile choices; they are uniform in the band, label,
rectangle copy, and point. The full definitions, including the support, compatibility, and
smooth-extension conditions, are Definitions~\ref{def:6.4} and~\ref{def:6.5}. Angular mean coefficients may still
depend on the auxiliary torus.

\begin{definition}[Regularity at the axis]\label{def:3.2}
A leading field is regular at the axis if its velocity
and pressure extend smoothly across \NSi{NS-F-P018-017}{r=0} in Cartesian coordinates, for
\NSi{NS-F-P018-018}{t<1}. For scalar
auxiliary profiles, ``regular at the axis'' means a smooth extension in
\NSi{NS-F-P018-019}{(X,\eta)} to \NSi{NS-F-P018-020}{X=0}.
\end{definition}

\begin{definition}[Parameter smallness and order of choices]\label{def:3.3}
For \NSi{NS-F-P018-021}{\alpha\geq0} and \NSi{NS-F-P018-022}{\beta>0}, we write
\NSi{NS-F-P018-023}{\alpha\ll\beta} when \NSi{NS-F-P018-024}{\alpha/\beta} is required to be sufficiently small, with the allowed size depending on
all data already fixed. In a chain of constants, choices proceed from right to left; small
reciprocals specify large parameters. For expressions depending on profile variables, this
smallness is uniform on the stated domain and is ensured by those parameter choices.
\end{definition}

The following table collects the principal symbols used in the construction, with references
to their definitions. The different background fields are grouped together. Repeated letters
are identified by their role or the section in which they are used. Orders in the coefficient
classes refer to normalized representatives, with the weights and derivative bounds specified
in their definitions.
